\section{System Model and Problem Formulation}
\label{sec:system-model}

In this section, we define the measurements a serving BS observes, set out what a window of them can add to a single report, and state the two estimation tasks.

\subsection{Measurement Model}
\label{subsec:measurement-model}

% src: FINAL_REPORT.md §3.1
At discrete snapshot time~$t$ the serving BS observes the measurement vector
\begin{equation}
  \mathbf{m}(t) = \big[\, \mathrm{RSS}(t),\ \mathrm{SINR}(t),\ \theta_{az}(t),\ \theta_{el}(t) \,\big],
  \label{eq:measurement-vector}
\end{equation}
with the azimuth/elevation pair $(\theta_{az}, \theta_{el})$ present only for UEs whose antenna count $N_{ant} \ge 2$, i.e. \emph{multi-antenna devices}, as opposed to \emph{single-antenna devices}. The channel experienced by multi-antenna UEs is modeled by maximum-ratio combining over the per-element channel,
\begin{equation}
  H_{\text{eff}}(f) = \sqrt{\textstyle\sum_{i=1}^{N_{ant}} \lvert H_i(f) \rvert^2},
  \label{eq:mrc}
\end{equation}
with a device-specific antenna-gain offset applied additively in dB.

% src: FINAL_REPORT.md §3.3 (window definition)
Given a history depth~$h$, the model input is the window of $h+1$ consecutive measurement vectors,
\begin{equation}
  \mathcal{W}_h(t) = \left[\, \mathbf{m}(t-h),\ \dots,\ \mathbf{m}(t) \,\right],
  \label{eq:transition-history}
\end{equation}
where $h=0$ is the single-report input. %UEs with $N_{ant} \ge 2$ are called multi-antenna devices, and those with $N_{ant} = 1$ single-antenna devices.

\subsection{What History Can Add}
\label{subsec:history-dof}
% src: FINAL_REPORT.md §1.1, §3.2; model: Italiano et al. 2025 Eqs. (14)-(15) (p. 14); TR 38.901 §7.4.1 (log-distance PL, log-normal shadow fading sigma_SF; itecspec.com, verified 2026-10-30)
Let $\mathbf{p}_{UE}$ and $\mathbf{p}_{BS}$ denote the UE and serving-BS positions, and let $r = \lVert \mathbf{p}_{UE} - \mathbf{p}_{BS} \rVert$ be their range. Under the log-distance path-loss model with log-normal shadow fading~\cite{italiano2025tutorial,3gpp38901},
\begin{equation}
  \mathrm{RSS} \approx P_{tx} - 10\gamma\log_{10} r + \chi,
  \label{eq:path-loss}
\end{equation}
where $P_{tx}$ is the transmit power, $\gamma$ the path-loss exponent and $\chi$ zero-mean shadow fading of standard deviation~$\sigma$. A single RSS value therefore identifies $r$ up to $\chi$ but leaves the bearing unconstrained, which is the distance-ring ambiguity of Section~\ref{subsec:ambiguity-hypothesis}.

A window can add to a single report in two distinct ways, which make different predictions and are tested separately in this work.

% src: FINAL_REPORT.md §3.3; csi_dataset.py build_derived_features (first differences of RSS, azimuth and ray coordinates divided by delta_t)
\noindent\textbf{Rates of change.} Consecutive reports of one UE, a time $\Delta t$ apart, give finite differences such as $[\mathrm{RSS}(t)-\mathrm{RSS}(t-1)]/\Delta t$, which approximate the time derivatives of the measurements; a window of depth $h$ gives $h$ such differences and therefore also their evolution. The raw window contains the consecutive values from which these differences follow, and the macro-cell features also supply the first differences of RSS, azimuth and the two ray coordinates explicitly~(Section~\ref{subsec:features}). Two of the derivatives carry geometric meaning. The rate of change of RSS,
\begin{equation}
  \frac{d\,\mathrm{RSS}}{dt} \approx -\frac{10\gamma}{\ln 10} \cdot \frac{\dot r}{r},
  \label{eq:drss-dt}
\end{equation}
reflects the radial component of the motion, and, where AoA is available, the angular rate,
\begin{equation}
  \dot\theta_{az} = \frac{v_\perp}{r},
  \label{eq:dtheta-dt}
\end{equation}
gives the tangential component $v_\perp$ of the velocity and, jointly with $\dot r$, the heading.

% src: FINAL_REPORT.md §7.5 (dRSS/dr expression); Equation (path-loss) above
Without AoA, only the first of these is available, and its usefulness is bounded by the path-loss gradient. Differentiating Equation~(\ref{eq:path-loss}) with respect to range gives
\begin{equation}
  \left\lvert \frac{d\,\mathrm{RSS}}{dr} \right\rvert = \frac{10\gamma}{r \ln 10}\ \text{dB/m},
  \label{eq:rss-gradient}
\end{equation}
so a step of fixed length moves RSS by an amount that falls as $1/r$, while the shadow-fading spread $\sigma$ does not shrink with range. RSS differences are therefore expected to carry less information far from the BS; Section~\ref{sec:results-no-gain} tests this expectation. Because a difference depends on which report precedes which, this mechanism requires the order of the reports.

% src: FINAL_REPORT.md §3.3; technical_documentation_revised.md 6.5 (averaging account of the angle history)
\noindent\textbf{Repeated noisy samples.} Over a short window the UE moves little compared with its range, so the reports are also several noisy samples of nearly the same quantities. Averaging them reduces the measurement error, such as the AoA estimation noise and quantization, independently of any information about the motion. This mechanism requires neither the order of the reports nor any motion, and it adds nothing when the measurements are exact.

The two mechanisms are therefore distinguished by whether the order of the reports matters, whether an unordered mean of the window recovers the gain, and whether the gain survives exact measurements; Sections~\ref{sec:results-window} and~\ref{sec:results-mechanism} apply these tests.

\subsection{Task Formulations}
\label{subsec:task-formulations}

% src: FINAL_REPORT.md §3.4; technical_documentation.md §7.4 (625-class memory limit)
We evaluate two task formulations. \emph{Classification} predicts one of $G$ discrete grid cells and is appropriate when the grid spacing greatly exceeds the position jitter. \emph{Regression} predicts a continuous position relative to the serving BS and is scored by the 2D Euclidean error. In the two larger deployments, UEs still move between the points of a $25\times25$ grid, but a classifier over its $625$ points is memory-prohibitive, so position is estimated by regression instead.
